\scanpage{90}
\mspage{18}
\uuline{Dictionnaire}\qquad Fonctions L \textendash\ \unsure{Chronologie|Cohomologie|Chronologique} à action galoisienne.

Soit $k$ un \deleted{corps ou} anneau

\deleted{\unsure{\(M\)|\(A\)} la catégorie} \deleted{$\operatorname{Mod}(\mathbb{Z})$ la catégorie} $C$ la catégorie des modules de longueur finie sur $k[t]$ $=$ \deleted{\unsure{vectoriels}}\inserted{modules} $V$ \inserted{proj.} de \deleted{dim}\inserted{rang} fini sur $k$ muni d'un endomorphisme $f$ $=$ $k[\deleted{\unsure{\mathbb{Z}}}\mathbb{N}]$-modules, \inserted{projectifs} de \deleted{\unsure{rang|de type}} finie sur $\deleted{\unsure{V}}k$.

Cette catégorie donne lieu à un $\lambda$-anneau, noté $R^{\unsure{\bullet|\cdot}}(\deleted{\unsure{\mathbb{Z}}}\mathbb{N};k)$, qui est aussi, \inserted{si $k$ est un corps,} le \deleted{$\lambda$-anneau} \inserted{groupe} libre engendré par les classes des \deleted{\unsure{\(V\)}}$k[\deleted{\unsure{\mathbb{Z}}}\mathbb{N}]$-modules simples \inserted{finis sur $k$}, ou avec les $k[t]$-modules simples, ou avec les polynômes \inserted{unitaires} irréductibles $P\in k[t]$, en associant à $P$ le module $k[t]/Pk[t]$. À tout tel $P$ \inserted{$k$ \deleted{\unsure{un}} corps ou non}, \inserted{est} associé à un $(V,f)$, dont le polynôme car. est précisément $P$,
\[
P(t)=\det(t\,\mathrm{id}-f)
\]
[vrai pour tt $P$, irréd. ou non !]. Il nous sera plus commode de remplacer $P(t)$ par
\[
Q(t)=t^dP\left(\frac{1}{t}\right)=\det(1-tf)
\]
qui est un polynôme quelconque à coeff. cst égal : $1$.
% archive sheet number: 89
